\documentclass[10pt]{article}
\usepackage[margin=1in]{geometry}
\usepackage{booktabs}
\usepackage{hyperref}
\usepackage{enumitem}
\usepackage{parskip}
\usepackage{microtype}
\usepackage{amsmath}
\usepackage{titlesec}

\titleformat{\section}{\normalfont\bfseries}{\thesection.}{0.5em}{}
\titlespacing{\section}{0pt}{1.2ex}{0.5ex}

\hypersetup{colorlinks=true, linkcolor=blue, urlcolor=blue, citecolor=blue}
\newcommand{\supp}{\href{[ANONYMISED\_REPO\_URL]}{supplementary repository}}

\begin{document}

\begin{center}
  {\large\bfseries Rebuttal --- Reviewer 4}
\end{center}

We thank the reviewer for a challenging critique identifying genuine gaps in
how the paper frames its contributions. We address each honestly.

% ---------------------------------------------------------------
\section*{Formal Incentive Compatibility}

\textit{``The mechanism does not have rigorous truthtelling guarantees\ldots the
truthfulness framing reads more like an analogy than any rigorous definition.''}

The reviewer is correct. LMSR is strictly proper when resolution is against
external ground truth; under consensus resolution that guarantee does not apply,
and the phrase ``our claims concern long-run coordination and robustness
properties'' is too vague as a substitute for a formal argument. We concede
there is no formal truthtelling guarantee in our setup.

What we can claim rigorously is more modest: the LMSR cost function creates
asymmetric costs between truth-seeking and deceptive behaviour when agents are
imperfectly adversarial, and these costs compound over episodes into systematic
wealth transfer. This is an empirical property of the mechanism interacting with
observed LLM behaviour, not a theoretical guarantee. We will reframe the paper's
incentive claims accordingly.

% ---------------------------------------------------------------
\section*{Prompts Do Not Mention Wealth ---
``Is this just debate with a scoring formula?''}

\textit{``The prompt provided on page 13 nowhere mentions wealth\ldots it's
unclear if it demonstrates that the system is mainly a debate-like framework
where the market is just an external scoring formula.''}

The reviewer correctly observes that agents are not told their capital. This
is intentional: wealth redistribution operates at the system level via the LMSR
cost function, regardless of agent awareness. Agents only need some correlation
with truth; the market amplifies that signal over episodes.

To test whether this design choice is consequential, we ran an ablation
augmenting each agent's prompt with its current capital balance and the statement
that wealth determines future influence (\supp, Section~5; Llama~3.1~8B,
TruthfulQA, 2v3 Informed, 3 trials):

\begin{center}\small
\begin{tabular}{lcc}
\toprule
\textbf{Condition} & \textbf{Market Acc} & \textbf{Mal. Wealth Share} \\
\midrule
No wealth in prompt (paper setup) & 73.3\% $\pm$ 7.2\% & 33.3\% \\
Wealth in prompt                  & 61.1\% $\pm$ 3.1\% & 67.2\% \\
\bottomrule
\end{tabular}
\end{center}

When adversarial agents know their capital determines future influence, accuracy
drops $12.2$pp and malicious wealth capture nearly doubles. This confirms the
mechanism is not purely instruction-following: if it were, wealth information
would have no effect. The no-wealth design deliberately prevents a class of
exploitation that wealth-aware adversaries execute.

We are not claiming LLMs explicitly optimise for capital. The LMSR creates
asymmetric outcomes: a lie requires a capital commitment forfeited at resolution,
unlike debate where persuasion costs nothing. This asymmetry compounds regardless
of agent awareness.

% ---------------------------------------------------------------
\section*{Experimental Assumptions and Scope}

\textbf{Strategic adversaries.}
The reviewer is correct that a rational, capital-maximising coalition could
dominate by outspending honest agents; our adversaries follow fixed directional
prompts, not wealth-optimisation. We should have stated this assumption
explicitly. The wealth ablation shows even partial wealth-awareness makes
adversaries more effective; fully strategic adversaries remain an open problem
we will flag in the limitations.

\textbf{Dataset memorisation.}
GPQA Diamond is a graduate-level benchmark substantially newer than most
pre-training data; market accuracy ($39$--$65\%$ under adversarial attack
vs.\ $25\%$ random baseline) is consistent with TruthfulQA, suggesting genuine
signal beyond memorisation. Truly out-of-distribution evaluation would
strengthen the claim and we will note this in Section~7.

\textbf{MCQ degradation.}
On GPQA MCQ (2v3 Informed), market accuracy falls to $21.3\%$---below the $25\%$
random baseline. The mechanism fails here and we do not dispute it. We will
revise the abstract and contributions to scope robustness claims to binary
settings upfront.

% ---------------------------------------------------------------
\section*{Revision Commitments}

\begin{itemize}[noitemsep,topsep=2pt,leftmargin=*]
  \item Reframe incentive claims as empirical; replace ``truthtelling incentives''
        with ``asymmetric cost structure that amplifies calibration advantage''
  \item State in Appendix~B.2 that agents do not observe their capital, and why
  \item Include wealth-prompt ablation as evidence against pure instruction-following
  \item State the imperfect-adversary assumption in the experimental setup;
        flag strategic adversaries as future work
  \item Scope robustness claims to binary settings in abstract and contributions;
        add dataset memorisation caveat to Section~7
\end{itemize}

\end{document}
